\documentclass[12pt,letterpaper]{article}
% -------------------------------------------------
% Formato general tipo APA 7.ª edición
% -------------------------------------------------
\usepackage[utf8]{inputenc}
\usepackage[T1]{fontenc}
\usepackage[spanish,es-tabla]{babel}
\shorthandoff{"} % desactiva los shorthands de comillas dobles de babel (" activa) para listings
\usepackage{geometry}
\usepackage{setspace}
\usepackage{indentfirst}
\usepackage{mathptmx}
\usepackage[table]{xcolor}
\usepackage{hyperref}
\usepackage{xurl} % mejora los cortes de línea en \url/\nolinkurl (útil dentro de tablas)
\usepackage{longtable}
\usepackage{array}
\usepackage{booktabs}
\usepackage{caption}
\usepackage{needspace}
\usepackage{enumitem}
\usepackage{fancyhdr}
\usepackage{titlesec}
\usepackage{listings}
\usepackage[most]{tcolorbox}
\usepackage{graphicx}
\usepackage{ragged2e}
\usepackage{float}
\usepackage{etoolbox}
% -------------------------------------------------
% Mejora visual de tablas
% -------------------------------------------------
\definecolor{tableStripe}{HTML}{F2F4F7}
\renewcommand{\arraystretch}{1.15}
\setlength{\arrayrulewidth}{0.8pt}
\setlength{\tabcolsep}{8pt}
\setlength{\LTpre}{0pt}
\setlength{\LTpost}{0pt}
\captionsetup[table]{font=small,labelfont=bf}
\AtBeginEnvironment{longtable}{%
  \setlength{\tabcolsep}{5pt}%
  \rowcolors{2}{tableStripe}{white}%
}
% Márgenes APA: 1 pulgada
\geometry{margin=1in}
\setlength{\headheight}{15pt}
% Interlineado APA: doble espacio
\doublespacing
% Sangría APA: primera línea de cada párrafo 0.5 pulgadas
\setlength{\parindent}{0.5in}
\setlength{\parskip}{0pt}
% Listas con sangría controlada
\setlist[itemize]{leftmargin=0.5in}
\setlist[enumerate]{leftmargin=0.5in}
% Alineación recomendada por APA
\raggedright
% -------------------------------------------------
% Encabezado y numeración APA
% -------------------------------------------------
\pagestyle{fancy}
\fancyhf{}
\rhead{\thepage}
\renewcommand{\headrulewidth}{0pt}
% -------------------------------------------------
% Colores y enlaces
% -------------------------------------------------
\definecolor{apaBlue}{HTML}{1F4E79}
\definecolor{codeBg}{HTML}{F6F8FA}
\definecolor{codeFrame}{HTML}{D0D7DE}
\definecolor{quoteBg}{HTML}{F7F7F7}
\definecolor{quoteFrame}{HTML}{BDBDBD}
\hypersetup{
    colorlinks=true,
    linkcolor=black,
    urlcolor=apaBlue,
    citecolor=black,
    pdftitle={Manual de Código: Jeiggar Vacation},
    pdfauthor={NICK ORTEGA}
}
% -------------------------------------------------
% Jerarquía de títulos APA
% -------------------------------------------------
\titleformat{\section}
  {\normalfont\centering\bfseries\Large}
  {\thesection.}{0.5em}{}
\titleformat{\subsection}
  {\normalfont\bfseries\large}
  {\thesubsection}{0.5em}{}
\titleformat{\subsubsection}
  {\normalfont\bfseries\itshape\normalsize}
  {\thesubsubsection}{0.5em}{}
\titlespacing*{\section}{0pt}{1.5em}{1em}
\titlespacing*{\subsection}{0pt}{1.2em}{0.8em}
\titlespacing*{\subsubsection}{0pt}{1em}{0.6em}
% -------------------------------------------------
% Código fuente
% -------------------------------------------------
\lstdefinelanguage{typescript}{
    keywords={const,let,var,function,return,export,import,from,type,interface,extends,if,else,await,async,null,true,false,default,new,class},
    sensitive=true,
    comment=[l]{//},
    morecomment=[s]{/*}{*/},
    morestring=[b]',
    morestring=[b]",
    morestring=[b]`
}
\lstdefinestyle{codeblock}{
    backgroundcolor=\color{codeBg},
    frame=single,
    rulecolor=\color{codeFrame},
    basicstyle=\ttfamily\small,
    breaklines=true,
    breakatwhitespace=false,
    columns=fullflexible,
    keepspaces=true,
    showstringspaces=false,
    % Bloques largos SÍ pueden cruzar páginas (no samepage).
    % Para bloques cortos que NO deben cortarse se usa \needspace antes.
    literate={á}{{\'a}}1 {é}{{\'e}}1 {í}{{\'i}}1 {ó}{{\'o}}1 {ú}{{\'u}}1
             {Á}{{\'A}}1 {É}{{\'E}}1 {Í}{{\'I}}1 {Ó}{{\'O}}1 {Ú}{{\'U}}1
             {ñ}{{\~n}}1 {Ñ}{{\~N}}1 {ü}{{\"u}}1 {Ü}{{\"U}}1
             {→}{{$\to$}}1 {≥}{{$\geq$}}1 {≤}{{$\leq$}}1 {·}{{$\cdot$}}1
             {*}{{*}}1 {!}{{!}}1 {i}{{i}}1 {OK}{{OK}}1 {X}{{X}}1,
    tabsize=2
}
\lstset{style=codeblock}
\newtcolorbox{infobox}{
    colback=quoteBg,
    colframe=quoteFrame,
    boxrule=0.8pt,
    arc=1mm,
    left=2mm, right=2mm,
    top=1mm,  bottom=1mm
}
\newcommand{\code}[1]{\nolinkurl{#1}}
\setlength{\LTpre}{0.8em}
\setlength{\LTpost}{0.8em}

% -------------------------------------------------
% Macros para diagramas
% -------------------------------------------------
% [H]: posición fija exacta donde se invoca (paquete float).
% Usado para TODOS los diagramas para evitar que aparezcan
% lejos de su contexto.
\newcommand{\diagramfigurehere}[3]{%
\begin{figure}[H]
\centering
\IfFileExists{#1}{%
    \includegraphics[width=0.9\textwidth,height=0.65\textheight,keepaspectratio]{#1}%
}{%
    \fbox{%
        \begin{minipage}[c][2.5in][c]{0.88\textwidth}
        \centering
        \textbf{Espacio reservado para el diagrama}\\[0.5em]
        Inserta aquí la imagen correspondiente en la ruta:\\[0.3em]
        \texttt{#1}
        \end{minipage}%
    }%
}
\caption{#2}
\label{#3}
\end{figure}
}

% Alias semántico para el diagrama ER de página completa
\newcommand{\diagramfigurefullpage}[3]{%
\begin{figure}[H]
\centering
\IfFileExists{#1}{%
    \includegraphics[width=\textwidth,height=0.85\textheight,keepaspectratio]{#1}%
}{%
    \fbox{%
        \begin{minipage}[c][4in][c]{0.95\textwidth}
        \centering
        \textbf{Espacio reservado para el diagrama}\\[0.5em]
        Inserta aquí la imagen correspondiente en la ruta:\\[0.3em]
        \texttt{#1}
        \end{minipage}%
    }%
}
\caption{#2}
\label{#3}
\end{figure}
}

% Compatibilidad: \diagramfiguresmall ahora es igual a \diagramfigurehere
\newcommand{\diagramfiguresmall}[3]{\diagramfigurehere{#1}{#2}{#3}}

\begin{document}
% -------------------------------------------------
% Portada tipo APA 7 estudiante
% -------------------------------------------------
\begin{titlepage}
    \centering
    \vspace*{1.5in}
    {\bfseries\Large Manual de Código: Plataforma Web de Gestión de Destinos y Administración de Servicios Turísticos\par}
    \vspace{0.5in}
    {Estudiantes de IV Semestre\par}
    \vspace{0.2in}
    {Fundación de Estudios Superiores Comfanorte (FESC)\par}
    {Ingeniería de Software\par}
    \vspace{0.2in}
    {Manual de Codigo del Proyecto\par}
    \vspace{0.2in}
    {Docente: Norly Alexandra Aguilar Quintero\par}
    \vspace{0.2in}
    {Mayo de 2026\par}
\end{titlepage}
\newpage
\tableofcontents
\newpage
\listoffigures
\newpage

% =================================================
\section{Visión General}
% =================================================

\textbf{Jeiggar Vacation} es una plataforma web de agencia de viajes turísticos, pensada para el mercado colombiano con proyección internacional. Permite a los usuarios explorar destinos nacionales e internacionales, cotizar planes por WhatsApp y acceder a información de servicios complementarios como asistencia médica, transporte, conectividad y trámites.

El sistema se divide en dos grandes áreas:

\Needspace{7\baselineskip}
\begin{longtable}{|>{\bfseries}p{0.25\textwidth}|p{0.65\textwidth}|}
\hline
Área & Descripción \\
\hline
Web Pública & Catálogo de destinos, mapas interactivos, formularios de cotización y páginas institucionales. \\
\hline
Panel Admin & CMS protegido por autenticación para gestionar destinos mediante CRUD completo con subida de imágenes. \\
\hline
\end{longtable}

\subsection*{Flujo de datos de alto nivel}

\diagramfigurehere{diagramas/flujo-datos-alto-nivel.png}{Flujo de datos de alto nivel del sistema Jeiggar Vacation.}{fig:flujo-datos-alto-nivel}

% =================================================
\section{Stack Tecnológico}
% =================================================

\subsection{Dependencias de producción}

\Needspace{16\baselineskip}
\begin{longtable}{|p{0.32\textwidth}|p{0.18\textwidth}|p{0.40\textwidth}|}
\hline
\textbf{Paquete} & \textbf{Versión} & \textbf{Uso} \\
\hline
\code{react} & \code{\string^19.2} & UI reactiva \\
\hline
\code{react-dom} & \code{\string^19.2} & Renderizado DOM \\
\hline
\code{react-router-dom} & \code{\string^7.13} & Enrutamiento SPA \\
\hline
\code{@supabase/supabase-js} & \code{\string^2.103} & Cliente BaaS: DB, Auth y Storage \\
\hline
\code{tailwindcss} & \code{\string^4.2} & Utilidades CSS \\
\hline
\code{@tailwindcss/vite} & \code{\string^4.2} & Plugin Vite para TailwindCSS 4 \\
\hline
\code{lucide-react} & \code{\string^1.3} & Iconografía SVG \\
\hline
\code{maplibre-gl} & \code{\string^5.21} & Mapas interactivos para destinos nacionales \\
\hline
\code{radix-ui} & \code{\string^1.4} & Primitivos accesibles de UI \\
\hline
\code{shadcn} & \code{\string^4.0} & Componentes UI sobre Radix \\
\hline
\code{aos} & \code{\string^2.3} & Animaciones scroll \\
\hline
\code{@fontsource-variable/geist} & \code{\string^5.2} & Tipografía Geist \\
\hline
\code{clsx} + \code{tailwind-merge} & latest & Composición dinámica de clases \\
\hline
\end{longtable}

\subsection{Herramientas de desarrollo}

\Needspace{9\baselineskip}
\begin{longtable}{|p{0.32\textwidth}|p{0.18\textwidth}|p{0.40\textwidth}|}
\hline
\textbf{Paquete} & \textbf{Versión} & \textbf{Uso} \\
\hline
\code{vite} & \code{\string^7.3} & Bundler / dev server \\
\hline
\code{typescript} & \code{\string~5.9} & Tipado estático \\
\hline
\code{eslint} & \code{\string^9.39} & Linting \\
\hline
\code{@vitejs/plugin-react} & \code{\string^5.1} & Soporte JSX con Fast Refresh \\
\hline
\end{longtable}

% =================================================
\section{Arquitectura del Proyecto}
% =================================================

El proyecto sigue una \textbf{arquitectura Feature-Sliced} adaptada. Cada feature es un módulo autónomo con sus propias páginas, componentes y servicios.

\subsection{Estructura de carpetas}

% Este bloque es largo (~35 líneas) y puede cruzar páginas normalmente.
% \needspace garantiza que al menos las primeras líneas no queden solas.
\Needspace{8\baselineskip}
\begin{lstlisting}[language={}]
jeiggar-vacation/
|---- src/
|   |---- app/                    # Bootstrap: router, shell, scroll
|   |   |---- router.tsx          # Definicion centralizada de rutas
|   |   |---- AppShell.tsx        # Layout publico (header + footer + <Outlet/>)
|   |   `---- ScrollToTop.tsx     # Restaura scroll en cada navegacion
|   |
|   |---- auth/                   # Modulo de autenticacion
|   |   |---- AuthContext.tsx     # Provider: sesion, isAdmin, signIn/signOut
|   |   |---- authContext.ts      # createContext (separado para evitar ciclos)
|   |   |---- useAuth.ts          # Hook consumidor del contexto
|   |   `---- ProtectedRoute.tsx  # Guard: redirige a /admin/login si no autenticado
|   |
|   |---- domain/
|   |   `---- types/              # Tipos TypeScript del dominio
|   |       |---- Map.ts          # Destination, City, Country...
|   |       |---- Plans.ts        # Plan, PlanCard...
|   |       |---- Landing.ts      # Hero, Stat...
|   |       `---- ...             # Un archivo por entidad/seccion
|   |
|   |---- lib/                    # Utilidades compartidas
|   |   |---- supabase.ts         # Singleton del cliente Supabase
|   |   |---- storage.ts          # Helper para URLs publicas de Storage
|   |   |---- utils.ts            # cn() y helpers genericos
|   |   `---- whatsapp.ts         # Numero y generador de links de WhatsApp
|   |
|   |---- components/             # Componentes UI reutilizables
|   |   |---- ui/                 # Atomos: ImageCard, Badge, Spinner...
|   |   |---- forms/              # Componentes de formulario genericos
|   |   `---- common/             # Breadcrumbs, Navbar, Footer...
|   |
|   |---- hooks/                  # Custom hooks globales
|   |---- mocks/                  # Datos estaticos de relleno
|   `---- features/               # Modulos de negocio
|       |---- admin/
|       |---- contact/
|       |---- destinations/
|       |---- landing/
|       |---- map/
|       |---- quotes/
|       `---- services/
|
|---- public/                     # Assets estaticos publicos
|---- docs/                       # Documentacion adicional
|---- .env.local                  # Variables de entorno locales
|---- .env.example                # Plantilla de variables de entorno
|---- vite.config.ts
|---- tsconfig.app.json
`---- package.json
\end{lstlisting}

\subsection{Estructura interna de un Feature}

Cada feature sigue esta convención interna:

\Needspace{10\baselineskip}
\begin{lstlisting}[language={}]
features/
`---- <nombre-feature>/
    |---- pages/              # Paginas conectadas al router
    |---- components/         # Componentes especificos del feature
    |---- services/           # Llamadas a Supabase/API
    |---- hooks/              # Hooks propios del feature
    `---- index.ts            # Barrel export opcional
\end{lstlisting}

\subsection{Diagrama de módulos y dependencias}

\diagramfigurehere{diagramas/modulos-y-dependencias.png}{Diagrama de módulos y dependencias del proyecto.}{fig:modulos-dependencias}

\subsection{Sistema de rutas}

Las rutas se definen en \code{src/app/router.tsx} usando \textbf{React Router v7} con \code{createBrowserRouter}. Todas las páginas usan \code{React.lazy()} para code splitting automático.

\Needspace{17\baselineskip}
\begin{longtable}{|p{0.32\textwidth}|p{0.38\textwidth}|p{0.18\textwidth}|}
\hline
\textbf{Ruta} & \textbf{Componente} & \textbf{Protegida} \\
\hline
\code{/} & \code{Home} & No \\
\hline
\code{/nosotros} & \code{AboutUs} & No \\
\hline
\code{/destinos} & \code{DestinationsSelector} & No \\
\hline
\code{/destinos/nacionales} & \code{MapRoute} & No \\
\hline
\code{/destinos/internacionales} & \code{InternationalDestinations} & No \\
\hline
\code{/destinos/:slug} & \code{DestinationDetail} & No \\
\hline
\code{/servicios} & \code{ServicesPage} & No \\
\hline
\code{/cotizar} & \code{QuotePage} & No \\
\hline
\code{/admin} & \code{AdminDashboard} & Sí \\
\hline
\code{/admin/destinos} & \code{DestinationList} & Sí \\
\hline
\code{/admin/destinos/nuevo} & \code{DestinationForm} & Sí \\
\hline
\code{/admin/destinos/:id/editar} & \code{DestinationForm} & Sí \\
\hline
\end{longtable}

% =================================================
\section{Guía de Configuración}
% =================================================

\subsection{Prerrequisitos}

\begin{itemize}
    \item \textbf{Node.js} $\geq$ 20.x, recomendado LTS.
    \item \textbf{npm} $\geq$ 10.x.
    \item Una cuenta y proyecto activo en \href{https://supabase.com}{Supabase}.
\end{itemize}

\subsection{Instalación paso a paso}

\Needspace{12\baselineskip}
\begin{lstlisting}[language=bash]
# 1. Clonar el repositorio
git clone <url-del-repo>
cd jeiggar-vacation
# 2. Instalar dependencias
npm install
# 3. Configurar variables de entorno
cp .env.example .env.local
# Editar .env.local con tus credenciales reales
# 4. Levantar servidor de desarrollo
npm run dev
\end{lstlisting}

El servidor arrancará en \code{http://localhost:5173} por defecto.

\subsection{Variables de entorno}

Edita el archivo \code{.env.local} con los valores de tu proyecto Supabase: Panel $\rightarrow$ Project Settings $\rightarrow$ API.

\Needspace{7\baselineskip}
\begin{lstlisting}[language={}]
# URL del proyecto Supabase
VITE_SUPABASE_URL=https://xxxxxxxxxxxxxxxx.supabase.co
# Anon key (clave publica, segura para el frontend)
VITE_SUPABASE_ANON_KEY=sb_publishable_xxxxxxxxxxxxxxxx
\end{lstlisting}

\begin{infobox}
\textbf{Importante:} Nunca subas \code{.env.local} al repositorio. Está incluido en \code{.gitignore}. Utiliza \code{.env.example} como referencia para otros desarrolladores.
\end{infobox}

\begin{infobox}
\textbf{Nota de seguridad:} Se usa la \code{anon key}, que es una clave pública. La seguridad de los datos se gestiona mediante \textbf{Row Level Security (RLS)} en Supabase, no ocultando la clave.
\end{infobox}

\subsection{Scripts disponibles}

\Needspace{9\baselineskip}
\begin{longtable}{|p{0.22\textwidth}|p{0.28\textwidth}|p{0.40\textwidth}|}
\hline
\textbf{Script} & \textbf{Comando} & \textbf{Descripción} \\
\hline
Desarrollo & \code{npm run dev} & Dev server con HMR en puerto 5173 \\
Build & \code{npm run build} & Compila TS + bundle de producción en \code{/dist} \\
Preview & \code{npm run preview} & Sirve el build de producción localmente \\
Linting & \code{npm run lint} & Ejecuta ESLint sobre todo el código \\
\hline
\end{longtable}

% =================================================
\section{Documentación de Lógica Core}
% =================================================

\subsection{Cliente Supabase: \code{src/lib/supabase.ts}}

El cliente es un \textbf{singleton} inicializado una sola vez. Si las variables de entorno no están definidas, lanza un error descriptivo en tiempo de inicio, evitando errores silenciosos en runtime.

\Needspace{5\baselineskip}
\begin{lstlisting}[language=typescript]
export const supabase = createClient(supabaseUrl, supabaseAnonKey)
\end{lstlisting}

\textbf{Regla:} Nunca crear una instancia nueva de \code{createClient} en otro archivo. Siempre importar \code{supabase} desde \code{@/lib/supabase}.

\subsection{Helper de Storage: \code{src/lib/storage.ts}}

Convierte rutas relativas almacenadas en la base de datos a URLs públicas completas de Supabase Storage.

\Needspace{5\baselineskip}
\begin{lstlisting}[language=typescript]
getPublicStorageUrl(path: string, bucket = 'destinations'): string
\end{lstlisting}

El helper maneja tres casos de forma retrocompatible:

\Needspace{8\baselineskip}
\begin{longtable}{|p{0.28\textwidth}|p{0.30\textwidth}|p{0.32\textwidth}|}
\hline
\textbf{Caso} & \textbf{Input} & \textbf{Output} \\
\hline
URL absoluta legado & \code{https://...} & La misma URL sin cambios \\
Asset local Vite & \code{/assets/foto.jpg} & La misma ruta local \\
Ruta relativa nuevo & \code{eje-cafetero-1234.jpg} & URL completa de Supabase Storage \\
\hline
\end{longtable}

\subsection{Servicio de Destinos Públicos: \code{destinations.service.ts}}

Ubicación: \code{src/features/destinations/services/destinations.service.ts}

Este servicio centraliza todas las consultas \textbf{de lectura} que usa la web pública. No modifica datos.

\subsubsection*{Flujo de datos}

\diagramfigurehere{diagramas/flujo-servicio-destinos-publicos.png}{Flujo de datos del servicio de destinos públicos.}{fig:flujo-servicio-destinos}

\subsubsection*{Funciones disponibles}

\textbf{\code{fetchDestinationBySlug(slug: string)}}

\begin{itemize}
    \item Carga un destino completo por su slug URL-friendly.
    \item Hace JOIN con \code{cities} o \code{countries} según el tipo.
    \item Retorna \code{null} si no existe, código \code{PGRST116}, o lanza error para otros casos.
    \item Filtro activo: solo devuelve destinos con \code{is\_active = true}.
\end{itemize}

\textbf{\code{fetchFeaturedDestinations()}}

\begin{itemize}
    \item Destinos marcados como \code{is\_featured = true}.
    \item Ordenados por \code{sort\_order\_home ASC} para controlar el orden desde el admin.
    \item Usado en la sección de destacados del Home.
\end{itemize}

\textbf{\code{fetchInternationalDestinations()}}

\begin{itemize}
    \item Todos los destinos con \code{destination\_type = 'international'}.
    \item Ordenados alfabéticamente por nombre.
    \item Hace JOIN con \code{countries}, no con \code{cities}.
\end{itemize}

\subsubsection*{Normalización de datos}

Un patrón consistente en este servicio es la normalización del JOIN anidado de Supabase, que puede devolver un objeto o un array según el contexto:

\Needspace{6\baselineskip}
\begin{lstlisting}[language=typescript]
// Supabase puede retornar city como objeto {} o array [{}]
const city = Array.isArray(d.city) ? d.city[0] : d.city
\end{lstlisting}

También convierte las claves de \code{snake\_case} de la DB a \code{camelCase} de TypeScript:

\Needspace{8\baselineskip}
\begin{longtable}{|p{0.42\textwidth}|p{0.42\textwidth}|}
\hline
\textbf{BD: snake\_case} & \textbf{TypeScript: camelCase} \\
\hline
\code{short\_description} & \code{shortDescription} \\
\code{image\_url} & \code{imageUrl} \\
\code{city\_id} & \code{cityId} \\
\hline
\end{longtable}

\subsection{Formulario de Destinos Admin: \code{DestinationForm.tsx}}

Ubicación: \code{src/features/admin/pages/DestinationForm.tsx}

Este es el componente más complejo del sistema. Sirve tanto para \textbf{crear} como para \textbf{editar} un destino, detectando el modo según la presencia del parámetro \code{:id} en la URL.

\subsubsection*{Estado del componente}

\diagramfigurehere{diagramas/estado-destination-form.png}{Diagrama de estados del componente DestinationForm.}{fig:estado-destination-form}

\subsubsection*{Variables de estado principales}

\Needspace{14\baselineskip}
\begin{longtable}{|p{0.26\textwidth}|p{0.24\textwidth}|p{0.40\textwidth}|}
\hline
\textbf{Estado} & \textbf{Tipo} & \textbf{Descripción} \\
\hline
\code{form} & \code{typeof EMPTY} & Todos los campos del destino \\
\code{cities / countries} & \code{City[] / Country[]} & Listas para los selectores \\
\code{loading} & \code{boolean} & Muestra spinner mientras carga en modo edición \\
\code{saving} & \code{boolean} & Deshabilita el botón submit durante guardado \\
\code{error} & \code{string | null} & Mensaje de error visible al usuario \\
\code{imageFile} & \code{File | null} & Imagen principal pendiente de subir \\
\code{previewUrl} & \code{string | null} & Blob URL para preview local de la imagen \\
\code{showNewCity} & \code{boolean} & Muestra u oculta el subformulario de nueva ciudad \\
\code{showNewCountry} & \code{boolean} & Muestra u oculta el subformulario de nuevo país \\
\hline
\end{longtable}

\subsubsection*{Lógica de \code{handleSubmit}}

El proceso de guardado sigue este orden estricto para evitar guardar referencias rotas:

\Needspace{10\baselineskip}
\begin{lstlisting}[language={}]
1. Subir imagen a Supabase Storage -> obtener nombre de archivo
2. Construir payload con image_url = nombre del archivo
3. Llamar createDestination() o updateDestination() segun modo
4. Si OK -> navegar a /admin/destinos
5. Si error de slug duplicado (codigo 23505) -> mostrar mensaje especifico
\end{lstlisting}

\subsubsection*{Subformularios inline: Ciudad / País}

El formulario permite crear una ciudad o país \textbf{sin salir de la pantalla}. El flujo es:

\Needspace{13\baselineskip}
\begin{lstlisting}[language={}]
Clic "+ Nueva ciudad"
  -> showNewCity = true
  -> Se muestra subformulario
  -> Usuario completa nombre + lat + lng
  -> handleCreateCity()
      -> Subir imagen de ciudad (opcional)
      -> createCity() en Supabase
      -> Agregar ciudad a la lista local (setCities)
      -> Seleccionar la nueva ciudad automaticamente (set('city_id', data.id))
      -> showNewCity = false
\end{lstlisting}

\subsubsection*{Generación automática de slug}

\Needspace{15\baselineskip}
\begin{lstlisting}[language=typescript]
function generateSlug(name: string) {
  return name
    .toLowerCase()
    .normalize('NFD')                    // Descomponer acentos
    .replace(/[\u0300-\u036f]/g, '')     // Eliminar diacriticos (a->a)
    .replace(/[^a-z0-9\s-]/g, '')       // Solo alfanumericos y guiones
    .trim()
    .replace(/\s+/g, '-')               // Espacios -> guiones
    .replace(/-+/g, '-')                // Multiples guiones -> uno
}
// "Canon del Chicamocha" -> "canon-del-chicamocha"
\end{lstlisting}

El slug se autogenera mientras el usuario escribe el nombre en modo creación, pero es editable manualmente.

\subsection{Autenticación: \code{src/auth/}}

El módulo de auth usa el patrón \textbf{Context + Provider} de React con Supabase Auth.

\diagramfigurehere{diagramas/flujo-autenticacion.png}{Flujo general del módulo de autenticación.}{fig:flujo-autenticacion}

\textbf{Lógica de \code{isAdmin}:} La autenticación de Supabase verifica la identidad del usuario, pero el rol \code{admin} se verifica en la tabla \code{profiles} del proyecto. Un usuario puede tener sesión activa pero no ser admin. \code{ProtectedRoute} verifica ambas condiciones.

\Needspace{10\baselineskip}
\begin{lstlisting}[language=typescript]
// src/auth/AuthContext.tsx
const { data } = await supabase
  .from('profiles')
  .select('role')
  .eq('id', sess.user.id)
  .single()
setIsAdmin(data?.role === 'admin')
\end{lstlisting}

% =================================================
\section{Estándares de Código}
% =================================================

\subsection{Convenciones de nombrado}

\Needspace{16\baselineskip}
\begin{longtable}{|p{0.30\textwidth}|p{0.30\textwidth}|p{0.30\textwidth}|}
\hline
\textbf{Elemento} & \textbf{Convención} & \textbf{Ejemplo} \\
\hline
Componentes React & \code{PascalCase} & \code{DestinationForm}, \code{ImageCard} \\
Funciones y variables & \code{camelCase} & \code{handleSubmit}, \code{fetchDestinations} \\
Tipos e Interfaces & \code{PascalCase} & \code{Destination}, \code{City} \\
Constantes globales & \code{UPPER\_SNAKE\_CASE} & \code{WHATSAPP\_NUMBER} \\
Archivos de componentes & \code{PascalCase.tsx} & \code{DestinationList.tsx} \\
Archivos de servicios & \code{kebab-case.service.ts} & \code{destinations.service.ts} \\
Archivos de hooks & \code{camelCase.ts} & \code{useAuth.ts} \\
Archivos de tipos & \code{PascalCase.ts} & \code{Map.ts}, \code{Plans.ts} \\
Rutas URL & \code{kebab-case} en español & \code{/destinos/internacionales} \\
Slugs de destinos & \code{kebab-case} sin acentos & \code{canon-del-chicamocha} \\
\hline
\end{longtable}

\subsection{Estructura de un componente funcional}

El orden de declaraciones dentro de un componente sigue esta secuencia:

\Needspace{22\baselineskip}
\begin{lstlisting}[language=typescript]
export default function MiComponente({ prop1, prop2 }: Props) {
  // 1. Hooks del router (useParams, useNavigate)
  const { id } = useParams()
  // 2. Estado local
  const [loading, setLoading] = useState(false)
  // 3. Contextos y hooks personalizados
  const { session } = useAuth()
  // 4. Efectos (useEffect)
  useEffect(() => { ... }, [id])
  // 5. Handlers y funciones derivadas
  function handleSubmit() { ... }
  // 6. Early returns (loading, error, empty states)
  if (loading) return <Spinner />
  // 7. JSX principal
  return (
    <div>...</div>
  )
}
\end{lstlisting}

\subsection{Manejo de tipos}

\begin{itemize}
    \item No usar \code{any}. Usar \code{unknown} cuando el tipo es incierto y aplicar narrowing cuando sea necesario.
    \item Los tipos del dominio se definen en \code{src/domain/types/} y se importan desde ahí.
    \item Los tipos locales de un componente, como interfaces de props o estados internos, se declaran en el mismo archivo.
    \item Para tipos de Supabase, tipar la respuesta usando el tipo propio de la función y no \code{as any}.
\end{itemize}

\Needspace{7\baselineskip}
\begin{lstlisting}[language=typescript]
// Correcto
const city = Array.isArray(data.city) ? data.city[0] : data.city
// Incorrecto
const city = (data as any).city
\end{lstlisting}

\subsection{Importaciones}

Usar el alias \code{@/} para rutas absolutas desde \code{src/}. Nunca usar rutas relativas que suban más de un nivel, como \code{../../}.

\Needspace{9\baselineskip}
\begin{lstlisting}[language=typescript]
// Correcto
import { supabase } from '@/lib/supabase'
import type { Destination } from '@/domain/types/Map'
// Incorrecto
import { supabase } from '../../../lib/supabase'
\end{lstlisting}

El alias \code{@/} está configurado en \code{vite.config.ts} y en \code{tsconfig.app.json}.

\subsection{Convenciones de Servicios}

\begin{itemize}
    \item Un archivo de servicio solo exporta funciones async, nunca estado.
    \item Cada función debe retornar el tipo exacto o lanzar el error para que el componente lo maneje.
    \item Separación estricta: \code{destinations.service.ts} solo lectura para web pública; \code{destinationsAdmin.service.ts} CRUD para panel admin.
\end{itemize}

% =================================================
\section{Flujo de Trabajo}
% =================================================

\subsection{Añadir un nuevo destino como Administrador}

\Needspace{22\baselineskip}
\begin{lstlisting}[language={}]
1. Ir a /admin/destinos
2. Clic en "Nuevo destino"
3. Seleccionar tipo: Nacional o Internacional
4. Completar informacion basica:
   - Nombre (el slug se genera automaticamente)
   - Categoria (ej: "Naturaleza", "Cultural")
   - Descripcion corta (<=160 caracteres)
   - Descripcion completa
5. Subir imagen (drag & drop o clic)
6. Ubicacion:
   - Nacional -> seleccionar ciudad (o crear una nueva inline)
   - Internacional -> seleccionar pais (o crear uno nuevo inline)
   - Ingresar coordenadas lat/lng
7. Detalles adicionales: clima, actividades separadas por coma
8. Publicacion:
   - Activar "Publicado" para que sea visible en la web
   - Activar "Destacado en el home" si debe aparecer en portada
9. Clic en "Crear destino"
\end{lstlisting}

\subsection{Añadir un nuevo Feature al código}

Sigue estos pasos para mantener la consistencia arquitectónica:

\Needspace{6\baselineskip}
\begin{lstlisting}[language=bash]
# Ejemplo: agregar un feature de "blog"
mkdir -p src/features/blog/pages
mkdir -p src/features/blog/services
\end{lstlisting}

\begin{enumerate}
    \item Crear el tipo en \code{src/domain/types/Blog.ts}.
    \item Crear el servicio en \code{src/features/blog/services/blog.service.ts}.
    \item Crear la página en \code{src/features/blog/pages/BlogPage.tsx}.
    \item Registrar la ruta en \code{src/app/router.tsx}.
\end{enumerate}

\Needspace{8\baselineskip}
\begin{lstlisting}[language=typescript]
// En router.tsx
const BlogPage = lazy(() => import('@/features/blog/pages/BlogPage'))
// Dentro del array de children del AppShell:
{ path: 'blog', element: <BlogPage /> }
\end{lstlisting}

Finalmente, añadir enlace en la navegación si corresponde.

\subsection{Modificar un campo de la base de datos}

Cuando se añade un campo nuevo a una tabla de Supabase:

\begin{enumerate}
    \item Ejecutar la migración en Supabase mediante SQL Editor.
    \item Actualizar el tipo en \code{src/domain/types/}.
    \item Actualizar la query en el servicio correspondiente añadiendo el campo al \code{select}.
    \item Actualizar el mapeo de \code{snake\_case} a \code{camelCase}, si aplica.
    \item Actualizar el formulario admin si el campo debe ser editable.
\end{enumerate}

\subsection{Git: Convención de ramas y commits}

\subsubsection*{Ramas}

\Needspace{9\baselineskip}
\begin{longtable}{|p{0.24\textwidth}|p{0.32\textwidth}|p{0.34\textwidth}|}
\hline
\textbf{Prefijo} & \textbf{Uso} & \textbf{Ejemplo} \\
\hline
\code{feature/} & Nueva funcionalidad & \code{feature/filtro-destinos} \\
\code{fix/} & Corrección de bug & \code{fix/slug-duplicado} \\
\code{chore/} & Mantenimiento, dependencias & \code{chore/actualizar-supabase} \\
\code{docs/} & Documentación & \code{docs/manual-codigo} \\
\hline
\end{longtable}

\subsubsection*{Commits: Conventional Commits}

\Needspace{8\baselineskip}
\begin{lstlisting}[language={}]
feat: anadir filtro por categoria en destinos
fix: corregir normalizacion de city en fetchDestinationBySlug
chore: actualizar @supabase/supabase-js a 2.103
docs: anadir seccion de auth al manual
\end{lstlisting}

% =================================================
\section{Base de Datos}
% =================================================

\subsection{Tablas principales}

\diagramfigurefullpage{diagramas/modelo-entidad-relacion.png}{Modelo entidad-relación de la base de datos del sistema.}{fig:modelo-er}

\subsection{Campo \code{destination\_type}}

\Needspace{8\baselineskip}
\begin{longtable}{|p{0.25\textwidth}|p{0.35\textwidth}|p{0.30\textwidth}|}
\hline
\textbf{Valor} & \textbf{Descripción} & \textbf{Relación geográfica} \\
\hline
\code{national} & Destino dentro de Colombia & Vinculado a una \code{city} \\
\code{international} & Destino fuera de Colombia & Vinculado a un \code{country} \\
\hline
\end{longtable}

\begin{infobox}
\textbf{Regla de integridad:} Al guardar, el formulario establece \code{city\_id = null} si el tipo es internacional, y \code{country\_id = null} si es nacional. Esto evita referencias cruzadas inconsistentes.
\end{infobox}

\subsection{Supabase Storage}

El bucket \code{destinations} almacena todas las imágenes de destinos y ciudades.

\subsubsection*{Convención de nombres de archivo}

\Needspace{8\baselineskip}
{\small
\begin{longtable}{|p{0.22\textwidth}|p{0.33\textwidth}|p{0.33\textwidth}|}
\hline
\textbf{Tipo} & \textbf{Formato} & \textbf{Ejemplo} \\
\hline
Imagen de destino & \code{\{slug\}-\{timestamp\}.\{ext\}} & \code{canon-del-chicamocha-1746301234.jpg} \\
Imagen de ciudad & \code{city-\{slug\}-\{timestamp\}.\{ext\}} & \code{city-san-gil-1746301234.webp} \\
\hline
\end{longtable}
}

La ruta relativa, por ejemplo \code{canon-del-chicamocha-1746301234.jpg}, es lo que se guarda en la columna \code{image\_url} de la base de datos. La función \code{getPublicStorageUrl()} construye la URL completa en tiempo de renderizado.

% =================================================
\section*{Referencias}
\addcontentsline{toc}{section}{Referencias}
% =================================================

Documentación oficial de React. (s. f.). \textit{React}. Recuperado el 3 de mayo de 2026, de \url{https://react.dev/}

Documentación oficial de Supabase. (s. f.). \textit{Supabase Docs}. Recuperado el 3 de mayo de 2026, de \url{https://supabase.com/docs}

Documentación oficial de Tailwind CSS. (s. f.). \textit{Tailwind CSS}. Recuperado el 3 de mayo de 2026, de \url{https://tailwindcss.com/docs}

Documentación oficial de Vite. (s. f.). \textit{Vite}. Recuperado el 3 de mayo de 2026, de \url{https://vite.dev/guide/}

\end{document}